So viel über \textsc{Kritter}s Gutachten oder vielmehr über den Teil davon, der mit dem Gegenstand meiner gegenwärtigen Kritik in unmittelbarem Zusammenhang steht. Auf den zweiten Teil des Gutachtens, der die durch die Waisenpensionen entstehende Vergrößerung der Ausgaben betrifft, werde ich in der zweiten Abteilung dieser Denkschrift zurückkommen. Dass in jenem ersten Teil von \textsc{Kritter}s Gutachten nur die Oberfläche des Gegenstands berührt ist, wird, meine ich, durch die vorstehenden Ausführungen hinreichend gezeigt sein. Überlegungen oder Vermutungen darüber, wie es dazu kam, dass \textsc{Kritter} ein so oberflächliches, nicht einmal mit seinen sonstigen öffentlichen Äußerungen übereinstimmendes Gutachten abgab, würden mich hier zu weit führen. Dennoch möchte ich glauben, dass, wenn man, statt sich darauf zu beschränken, zwei Einzelfragen zu stellen,
\begin{quote}
von denen zudem die eine im Sinne des Fragestellers überhaupt keine bestimmte Antwort zuließ, den p. \textsc{Kritter} in irgendeiner Form an den weiteren Beratungen hätte teilnehmen lassen oder ihn wenigstens unter vollständiger Mitteilung aller Sachverhältnisse über die Angemessenheit des Plans zu Rate gezogen hätte, den man auf sein Gutachten gründen zu können glaubte,
\end{quote}
\noindent die Frage der Regelung aller künftigen Pensionserhöhungen nicht so, wie geschehen, über das Knie gebrochen worden wäre\footnote{Die Zahl der im Lauf meines obigen Berichts angeführten oder angedeuteten Beispiele laxer Geschäftsbehandlung hätte leicht noch vergrößert werden können; ich hielt das jedoch für ebenso unnötig wie unerfreulich.}. Das ist aber unterblieben. \textsc{Kritter} erhielt aus der Witwenkasse für sein Gutachten ein Honorar von 4 Thl. 24 mgr., so wie P. im Ministerialreskript vom 20. November 1794 eine Belobigung seiner gut ausgearbeiteten Denkschrift erhielt.

In diesem Reskript wurden die Vorschläge genehmigt. Hinsichtlich der Erhöhung der Pensionen wurde jedoch bestimmt, dass, wenn wider Erwarten \textit{unglückliche Umstände} künftig eine Verminderung der Pensionen nötig machen sollten, die dann vorhandenen Witwen diese hinnehmen müssten. Auch wurde für die künftig nach der Progressionsregelung vorzunehmenden Pensionserhöhungen die jeweilige Ratifikation vorbehalten. \textsuperscript{.} Es scheint bemerkenswert, dass hier nur an eine Notwendigkeit gedacht wird, die aus unglücklichen Umständen entsteht, nicht aber an diejenige, die möglicherweise aus einer zu großen Witwenzahl hervorgehen könnte. Hat man einen solchen Fall für unmöglich gehalten, oder hat man die Regelung im selben Sinn wie \textsc{Sch.} aufgefasst und es als selbstverständlich betrachtet, dass in einem solchen Fall die Pension wieder herabgehen müsse? Bei der Ratifikation der ersten Erhöhung (1799 Mai 23) wurde derselbe Vorbehalt wiederholt. Dabei wurde aber nur allgemein von Umständen gesprochen, die eine erneute Verminderung nötig machen könnten, ohne sie als unglückliche zu bezeichnen. Bei den späteren Ratifikationen ist, soweit ich habe feststellen können, der Vorbehalt nicht wiederholt worden.

